NBO attains the declared scalar-interval tolerance in 36 of its 36 complete runs, including 36 at the first declared fitting stage. Its largest returned mean-regret upper bound is $1.47\times10^{-6}$. The shared conditional actor attains the target in 3 of 36 runs; all remaining outcomes stay unresolved rather than being replaced. The conventional regressions and structure-exploiting simulation procedures are retained on the same accuracy and cost basis.

Across the twelve dimension--query-volume cells, the lowest mean service cost among methods certified in all three streams belongs to Quadratic in 8 cells, Cached SAA in 1 cells, Radial basis in 2 cells, NBO in 1 cells. This is an observed comparison among the declared procedures on the reported machine, not a statistical assertion of minimum possible computation. Validation work is included in every entry and becomes a substantial part of the large-query bill. Consequently, reducing the query cost alone need not produce the same ranking as reducing the total work to the economic target.

The complete comparison uses 288.476 measured seconds in its reported service environment. The sum of all method service bills is 268.750 seconds, and the separately identified post-stop mechanism calculations use 16.337 seconds. Shared startup and remaining driver and individual-record overhead are part of the full comparison, not retrospectively assigned to a selected winner. Serialization of the final aggregate summary and publication generation are outside that clock.
